\section{Elección de líder: Algoritmo de Bully}

Cuando el monitor declara caído al líder, los nodos restantes tienen que
elegir uno nuevo. El código está en el módulo \texttt{bully} del crate
\texttt{tp2\_reservas}. Esta parte se corre con
\texttt{cargo run -- parte2-lider}; la elección se dispara solo si el nodo
que cayó era el líder, porque si cae un seguidor el liderazgo no cambia.

\begin{figure}[H]
\centering
\begin{tikzpicture}[x=3.2cm,y=0.8cm]
  % Nodos
  \node[caja] (n0) at (0,0) {Nodo 0};
  \node[caja] (n1) at (1,0) {Nodo 1};
  \node[caja] (n2) at (2,0) {Nodo 2 (caído)};

  % Líneas de vida (empiezan un poco más abajo de las cajas)
  \foreach \x in {0,1,2} \draw[dashed] (\x,-0.4) -- (\x,-6.5);
  
  % Flecha 1: Eleccion{0}
  \draw[flecha] (0,-1.2) -- node[above,font=\scriptsize]{\texttt{Eleccion\{0\}}} (1,-1.2);
  
  % Flecha 2: Respuesta{1}
  \draw[flecha] (1,-2.0) -- node[above,font=\scriptsize]{\texttt{Respuesta\{1\}}} (0,-2.0);
  
  % Flecha 3: Eleccion{1} (hacia el caído)
  \draw[flecha] (1,-2.8) -- node[above,font=\scriptsize]{\texttt{Eleccion\{1\}}} (2,-2.8);
  \node[font=\scriptsize,anchor=west] at (2.2,-2.8) {sin respuesta};
  
  % Flecha 4: Timeout (en el nodo 1)
  \draw[flecha,dashed] (1,-3.6) -- node[above,font=\scriptsize]{timeout} (1.8,-3.6);
  
  % Flecha 5 y 6: Coordinador{1}
  \draw[flecha] (1,-4.4) -- node[above,font=\scriptsize]{\texttt{Coordinador\{1\}}} (0,-4.4);
  \draw[flecha] (1,-5.2) -- node[above,font=\scriptsize]{\texttt{Coordinador\{1\}}} (2,-5.2);
  \node[font=\scriptsize,anchor=west] at (2.2,-5.2) {(descartado)};
\end{tikzpicture}
\caption{Secuencia del Bully cuando cae el líder (nodo 2) y queda el nodo 1
como el vivo de mayor ID.}
\label{fig:bully-secuencia}
\end{figure}

\subsection{Supuestos}

Los nodos tienen IDs únicos y ordenados. Se comunican por canales con
timeout: el mismo \texttt{recv\_timeout} del monitor sirve para no esperar
para siempre una respuesta. Los fallos son de tipo \emph{crash-stop}: un
nodo o está vivo o no contesta, no hay respuestas corruptas ni
bizantinas. Estas tres hipótesis son las del paper original de
García-Molina (1982) y son las que el simulador respeta: los nodos caídos
simplemente no reciben ni contestan mensajes.

\subsection{Los tres mensajes}

El algoritmo usa tres tipos de mensaje, definidos en el enum
\texttt{MsgBully}:

\begin{itemize}
    \item \texttt{Eleccion\{de\}}: ``quiero ser líder, ¿hay alguien con ID
    mayor que yo vivo?''. Se manda a todos los nodos con ID mayor.
    \item \texttt{Respuesta\{de\}}: ``yo (ID mayor) estoy vivo, retirate''.
    Es el OK del paper.
    \item \texttt{Coordinador\{lider\}}: ``soy el nuevo líder, aceptame''.
    Lo manda el ganador a todos los demás.
\end{itemize}

Un nodo caído no recibe nada y no contesta nada. Para los demás es
silencio, igual que en el monitor: no hay un mensaje de error, hay
ausencia de mensaje.

\subsection{El escenario del trabajo}

Nodos 0, 1 y 2. El líder es el 2, que es el de mayor ID. El monitor
detecta que el 2 cayó y avisa al nodo 0. La secuencia es la de la
Figura~\ref{fig:bully-secuencia}:

\begin{enumerate}
    \item El 0 inicia la elección: manda \texttt{Eleccion\{0\}} a los
    nodos con ID mayor. Solo al 1, porque el 2 está caído.
    \item El 1 recibe \texttt{Eleccion\{0\}}, contesta
    \texttt{Respuesta\{1\}} al 0 y arranca su propia elección: manda
    \texttt{Eleccion\{1\}} a los mayores que él. Solo al 2, que no
    contesta.
    \item El 0 recibe \texttt{Respuesta\{1\}} y se retira. Hay alguien
    mayor vivo. Pasa a esperar el anuncio del 1.
    \item El 1 hace timeout sin respuestas (el 2 no contesta) y se
    declara líder. Manda \texttt{Coordinador\{1\}} a todos.
    \item El 0 recibe \texttt{Coordinador\{1\}} y lo acepta. El 2 sigue
    caído y no participa.
\end{enumerate}

El resultado es el nodo 1, que es el vivo con mayor ID. Los mensajes
intercambiados son pocos: dos de elección, una respuesta y un anuncio. Si
el 1 también estuviera caído, el 0 haría timeout sin respuestas y se
declararía líder directamente.

\subsection{Complejidad y garantía}

En el peor caso, todos los nodos inician elección a la vez y cada uno
manda \texttt{Eleccion} a todos los mayores: son $O(n^2)$ mensajes. En el
mejor caso, el nodo que detecta la falla es el de mayor ID vivo y se
declara líder con $O(n)$ mensajes. La garantía es que gana el nodo vivo
con mayor ID, siempre que los timeouts sean mayores que la latencia de la
red. Si no lo son, un nodo puede declararse líder antes de que llegue la
respuesta del mayor, y el algoritmo converge igual pero con más
elecciones espurias.

\subsection{Limitaciones}

Bully no resuelve el \emph{split-brain}. Si el monitor declara caído al
líder por una falsa alarma y el líder sigue vivo pero lento, pueden
coexistir dos líderes temporales: el viejo, que cree que sigue siéndolo, y
el nuevo, elegido por los demás. El paper original asume que el monitor no
se equivoca; en la práctica se resuelve con quórum (el líder necesita
mayoría para decidir), leases (el liderazgo vence solo) o fencing tokens
(cada líder tiene un número de generación y los recursos rechazan
operaciones de generaciones viejas). Bully no usa ninguna de las tres: es
un algoritmo didáctico, no uno de producción. Raft y Paxos existen
justamente porque Bully no alcanza para sistemas reales.

La segunda limitación es que Bully depende de que los IDs sean conocidos y
estables. En un sistema con nodos que se unen y se van dinámicamente, el
ID deja de ser un buen criterio de liderazgo: el nodo con mayor ID puede
ser el más nuevo y el menos cargado, y elegirlo líder no es
necesariamente lo que conviene.

\subsection{Pruebas}

\texttt{cae\_el\_lider\_y\_gana\_el\_mayor\_vivo} corre la elección con
tres nodos, tira abajo al 2 y verifica que el líder que ven los vivos es
el 1.

\texttt{cinco\_nodos\_caen\_los\_dos\_mayores} extiende el escenario a
cinco nodos, tira abajo al 3 y al 4, y verifica que gana el 2.

\texttt{si\_el\_mayor\_esta\_vivo\_gana\_aunque\_inicie\_otro} corre con
los tres nodos vivos e inicia la elección desde el 0. El 2 responde y
gana: el algoritmo no depende de quién inicia, depende de quién tiene el
ID más alto.

Los tres tests verifican la misma propiedad desde ángulos distintos: gana
el vivo con mayor ID, sin importar quién disparó la elección ni cuántos
nodos cayeron.